\subsection{Integral cycles detecting the twisted norm}\label{addc:section}

We construct the cycles and the common family of evaluations used in
\eqref{sup:lengthcount}. Write \(A=\mathbf Z/q\), \(q=\ell^\nu\),
\(w=\ell^c\), \(\kappa=1+w\), \(a=\kappa^E\), \(M=2E+1<w\).
Here \(\nu\ge c\), and \(c\ge2\) if \(\ell=2\). The layer \(F/L'\)
is cyclic with generator \(\gamma\), and contains a primitive
\(q\)-th root \(\zeta\) on which \(\gamma\) acts by \(\kappa\).
In this cyclotomic layer \(\gamma\) and \(\kappa\bmod q\) have the
same order. Write \(C_q=\langle h\rangle\) for the group of order \(q\).
All bar chains below have coefficients in \(A\).

\begin{lemma}[One integral test for each residue field]\label{addc:residue}
Let \(k\) be a finite field containing a primitive root
\(\xi\) of order \(q\), and let \(v:k^\times\to A\) be a
homomorphism. There are homology evaluations \(\theta_V\) of degree
\(M\) on \(\mathrm{GL}(V)\) for finite-dimensional \(k\)-spaces \(V\).
They are invariant under change of basis and additive under direct
sum of representations. Put
\[
 z=\sum_{j_1,\dots,j_E\bmod q}
       [h^{j_1}|h|\cdots|h^{j_E}|h]\in C_{2E}(C_q;A).
\]
If \(h\) acts on \(V=\bigoplus_x V_x\) by \(\xi^x\) on \(V_x\),
and \(g=\bigoplus_x g_x\) preserves the summands, then the prism
of \(g\) on the image of \(z\) has value
\(\sum_x x^E v(\det g_x)\).
\end{lemma}
\begin{proof}
The first face in the boundary of \(z\) occurs \(q\) times. The
other faces cancel in consecutive pairs after translating a summation
index by one, including the last face. Thus \(z\) is a cycle.
Write \(|k^\times|=q m\) and choose a generator with \(m\)-th power
\(\xi\). Such a generator exists since units modulo \(qm\) surject
onto units modulo \(q\), by prime-power reduction and the Chinese
remainder theorem. With \(L(b)\) its integer discrete logarithm in
\([0,qm)\), put
\[
   u(b,d)=\frac{L(b)+L(d)-L(bd)}{qm}\pmod q .
\]
This symmetric two-cocycle has zero coboundary by cancellation of
the four logarithmic differences. On \(k^\times\) take the cup
product \(u^{\smile E}\smile v\). The prism of \(g\) inserts \(g\)
in every position with alternating signs. Terms in the first \(2E\)
positions cancel in pairs by symmetry of \(u\). The final insertion
contributes \(v(g)\) times the product of the \(E\) carries. For
\(h\mapsto\xi^x\), summing each carry over \(j\bmod q\) gives
\(q L(\xi^x)/(qm)\equiv x\bmod q\) by telescoping. This gives
\(x^E v(g)\) in dimension one.

For general \(V\), sum this cocycle over the lines of \(V\).
More explicitly, lift a bar simplex through the action on
\(\mathbf P(V)(k)\), once for each starting line. Choose one nonzero
vector \(v_x\) on each line. For arrows \(g_1,\ldots,g_M\) and initial
line \(x_0\), put \(x_i=g_i^{-1}x_{i-1}\) and read the scalar \(b_i\)
from \(g_i v_{x_i}=b_i v_{x_{i-1}}\). Evaluate the cocycle on these
scalars. Composition multiplies the scalars and
the arrows permute the starting lines, so the sum commutes with the
bar faces. Changing the vectors gives naturally isomorphic functors
from the action groupoid to \(k^\times\). The simplicial prism of
that natural isomorphism gives the same value on cycles. This also
proves invariance under change of basis. For a block sum \(V\oplus W\),
lines with nonzero first projection lie over \(\mathbf P(V)\) with
\(|k|^{\dim W}\) lifts per line. Normalize their first projections
by the chosen vectors of \(V\). Their scalars are the scalars in \(V\).
The other lines are precisely \(\mathbf P(W)\).
Since \(|k|\equiv1\bmod q\), the evaluation is additive.

On a scalar block the prism value is a homomorphism in \(g\).
Indeed, composing the prisms for two successive arrows gives the
prism for their product modulo the boundary over their two-simplex.
Here both arrows centralize the cycle. This homomorphism factors
through the determinant. In dimension at least two the elementary
matrices generate \(\mathrm{SL}(V)\), and each is a commutator
with a diagonal matrix since \(|k|>2\). Its value on
\(\operatorname{diag}(b,1,\dots,1)\) follows from additivity and the
one-dimensional computation. Applying additivity to all character
blocks gives the formula. In this application each block prism is
the image of the same cycle on \(C_q\times\mathbf Z\).
\end{proof}

\begin{lemma}[Isolation of one cyclic direction]\label{addc:filter}
Set \(d_t(h)=(h,h^t)\) for \(0\le t\le M\) and
\(\chi_t(j_1,j_2)=j_1+t j_2\) on \(J=A^2\).
There are \(n_t\in A\), with \(n_{t_0}\) a unit, and a chain \(B\)
on \(C_q^2\) with \(\partial B=\sum_t n_t(d_t)_*z\).
For this \(t_0\) there is an \emph{integer}-valued function \(\eta\)
on \(J\) such that, in the integral group ring of the additive group
of \(A\),
\[
 (\chi_t)_*\eta=0\quad(t\ne t_0),\qquad
 (\chi_{t_0})_*\eta=\sum_{i=0}^M p_i[\kappa^i],\qquad
 \sum_i p_i X^i=P(X)=\prod_{j=0}^{M-1}(X-\kappa^j).
\]
Here pushforward adds coefficients in each fibre and \([x]\) is a
group-ring basis element.
\end{lemma}
\begin{proof}
Tensor the two cyclic periodic free resolutions over \(\mathbf Z\)
to resolve the trivial module for \(C_q^2\). After tensoring with
trivial coefficients \(A\), the maps \(h-1\) and the norm both
vanish. Thus \(H_{2E}(C_q^2;A)\simeq A^{2E+1}\). The \(M+1\)
classes give an \(M\)-by-\((M+1)\) matrix. Its Smith form over
\(\mathbf Z/\ell^\nu\) has a free zero column. Transporting its
coordinate vector back gives a relation with a unit entry, hence
the bounding chain \(B\).

Let \(U=X^w-1\) and let \(T\) act on \(\mathbf Z[X]\) by
\(Tf(X)=f(X^\kappa)\). Then \(Tf\equiv f\bmod U\) and
\(TU=(1+U)^\kappa-1\equiv\kappa U\bmod U^2\). Consequently
\((T-\kappa^j)(U^j f)\) is divisible by \(U^{j+1}\). Applying
the \(M\) factors to \(X\) shows that \(\sum_i p_i X^{\kappa^i}\)
is divisible by \(U^M\) over \(\mathbf Z\). Substitute \(X=[1]\)
in the group ring. Write the resulting element as
\(([w]-[0])^M H\). For each \(t\ne t_0\) choose \(b_t\in A\) with
\((t_0-t)b_t=w\). This is possible since \(0<|t_0-t|\le M<w\)
and \(q\) is an \(\ell\)-power. Lift \(H\) by \(x\mapsto(x,0)\)
and multiply it by
\(\prod_{t\ne t_0}([(-t b_t,b_t)]-[(0,0)])\).
Pushing the product forward gives precisely the two asserted
formulas. The argument also covers \(c=\nu\).
\end{proof}

\begin{theorem}[Simultaneous cycles for the arithmetic detector]\label{addc:cycle}
Suppose that \(\alpha\in F^\times\) and an integral-coefficient
divisor \(\mathfrak d\) satisfy
\(\operatorname{div}\alpha=(\gamma^{-1}-a)\mathfrak d\), with
\((\gamma^{-1}\mathfrak d)(P)=\mathfrak d(\gamma P)\).
There are a matrix size \(N\), a cycle \(y_\alpha\) of degree \(M\)
on \(\mathrm{GL}_N(\mathcal O_{L'})\), and a finite set of excluded
primes. At every remaining prime \(P\) of \(F\) of relative residue
degree one, reduction followed by Lemma~\ref{addc:residue}, with
\(\xi=\bar\zeta\), gives
\[
 \Phi_P(y_\alpha)=
  \prod_{\substack{0\le j<M\\j\ne E}}(a-\kappa^j)\,
       v_P^\times\!\left(\overline{N_{\gamma,a}\alpha}\right).
\]
The evaluation family can be fixed independently of \(\alpha\) once
the indices \(v_P^\times\) have been chosen.
\end{theorem}
\begin{proof}
Put \(Q(X)=P(X)/(X-a)\) and \(\beta=Q(\gamma^{-1})\alpha\),
where an integral group-ring operator acts multiplicatively on
\(F^\times\). Then
\(\operatorname{div}\beta=P(\gamma^{-1})\mathfrak d\).
Let \(\mathfrak a\) be the fractional ideal of divisor \(\mathfrak d\).
On \(F^J\) consider the two \(\mathcal O_{L'}\)-lattices
\[
 \Lambda_0=\bigoplus_{j\in J}\mathcal O_F,\qquad
 \Lambda_1=\bigoplus_{j\in J}\mathfrak a^{\eta(j)}.
\]
The diagonal representation \(\rho\) of \(C_q^2\), with characters
\(j\) evaluated at \(\zeta\), preserves both. We construct
\(L'\)-linear isomorphisms \(g_t:\Lambda_0\to\Lambda_1\) commuting
with \(\rho d_t\). For \(t\ne t_0\), use the blocks
\(\chi_t(j)=x\). The determinant ideal of each target block is the
unit ideal by Lemma~\ref{addc:filter}. Over the Dedekind domain
\(\mathcal O_F\), a sum of \(s>0\) fractional ideals is isomorphic
to \(\mathcal O_F^{s-1}\) plus their product. Indeed, after scaling
two ideals to coprime integral representatives their sum map onto
\(\mathcal O_F\) splits with kernel isomorphic to their product.
Iteration proves the assertion. Thus one may choose the map on
each block with determinant \(1\), correcting its determinant by
a unit on a source coordinate.

At \(t_0\), on each coordinate with \(\chi_{t_0}(j)=\kappa^i\)
modulo \(q\), change coordinates by the \(L'\)-automorphism
\(\gamma^{-i}\) of \(F\). This is unambiguous by the equality of
orders in the layer. The character on these coordinates becomes \(1\)
in additive exponent notation, so \(h\) acts as \(\zeta\).
Keep the other coordinates unchanged. On the combined exponent-one
block the determinant divisor of the transformed target is
\(\sum_i p_i\gamma^{-i}\mathfrak d=\operatorname{div}\beta\).
Every other block has determinant divisor zero. The same ideal
argument gives a map of determinant \(\beta\) on the exponent-one
block and \(1\) elsewhere. This block is nonempty since \(j=(1,0)\)
occurs. Undo the coordinate change to obtain \(g_{t_0}\).

Write \(H=\operatorname{Aut}_{\mathcal O_{L'}}(\Lambda_0)\),
\(g_*=g_{t_0}\), \(\rho'=g_*^{-1}\rho g_*\), \(k_t=g_*^{-1}g_t\).
Both \(\rho,\rho'\) and all \(k_t\) take values in \(H\).
In the bar construction write \(P_g\) for the homotopy from the
identity to conjugation by \(g^{-1}\). Form the chain in \(H\)
\[
 y=n_{t_0}^{-1}\left(\rho'_* B-\rho_* B+
           \sum_t n_t P_{k_t}(\rho'd_t)_*z\right).
\]
Indeed \(k_t^{-1}(\rho'd_t)k_t=\rho d_t\), so the prism boundary
formula and \(\partial B=\sum_t n_t(d_t)_*z\) give \(\partial y=0\).
In the group containing \(H\) and the \(g_t\), this cycle is
homologous to \(n_{t_0}^{-1}\sum_t n_t P_{g_t}(\rho d_t)_*z\).
To see this, apply the prism identity to \(\rho_*B\) at \(g_*\).
Modulo a boundary it replaces \(\rho'_*B-\rho_*B\) by
\(P_{g_*}\rho_*\partial B\). On each summand the two consecutive
prisms are homologous to the prism of \(g_*k_t=g_t\), using the
boundary of the triangle for the two composable conjugations.
The finite projective \(\mathcal O_{L'}\)-module \(\Lambda_0\)
is a summand of a finite free module. Extending the \(H\)-action
by the identity on a projective complement gives the asserted cycle
\(y_\alpha\) on an integral general linear group.

Exclude ramification, primes dividing \(q\), and all conjugates of
the supports of \(\operatorname{div}\alpha\) and \(\mathfrak d\).
Also discard the finitely many places needed to make the lattices
equal and the displayed maps and their inverses integral. These
exclusions make all conjugates of \(\alpha,\beta\) units before
the residue test. At a remaining relative-degree-one place the Galois extension
splits over the residue field \(k\). The reduction of \(\Lambda_0\)
therefore has coordinates \((\sigma,j)\), with
\(\sigma\in\operatorname{Gal}(F/L')\). In those coordinates conjugating
an \(F\)-coordinate by a power of \(\gamma\) merely permutes its
embeddings. Write \(\sigma\zeta=\zeta^{k_\sigma}\).
The trivial complement contributes zero in positive degree.
By Lemma~\ref{addc:residue}, every \(t\ne t_0\) contributes zero
because its block determinants are \(1\). The remaining prism
contributes \(\sum_\sigma k_\sigma^E v_P^\times(\overline{\sigma\beta})\).
The unit coefficient \(n_{t_0}\) has already been divided out.
Replacing \(\alpha\) by \(\gamma^{-i}\alpha\) in this last weighted
sum for \(\alpha\) multiplies it by \(a^i\), by reindexing
\(\sigma\mapsto\sigma\gamma^i\). For \(\alpha\) the sum is
\(v_P^\times(\overline{N_{\gamma,a}\alpha})\).
Applying \(Q\) gives the claimed factor \(Q(a)\). All residue
evaluations use the single construction of Lemma~\ref{addc:residue}.
\end{proof}

In particular one can use the union of the excluded sets for all
representatives of the finite image \(I\) in \eqref{sup:lengthcount}.
Stabilization preserves the value and repeating a block \(b'\) times
multiplies it by \(b'\), by direct-sum additivity. The exact valuation of \(Q(a)\) is
\((M-1)c+2v_\ell(E!)\): the distances \(|j-E|\) run twice through
\(1,\ldots,E\), and \(v_\ell(\kappa^n-1)=c+v_\ell(n)\) by the
binomial formula, with \(c\ge2\) at two. Thus the scaled test
\(\ell^{M-1}\Phi\) already detects the norm raised to the power
\(\ell^{M-1+(M-1)c+2v_\ell(E!)}\), which suffices for the stated
choice of \(b\). The finite-direction filter, the residue cocycle
and the assembly on the two lattices have corresponding constructions
in the pinned source.
\src{src/Development/AutomorphicArithmetic/AutomorphicVarieties/ModularityArgument/ClassCounting/ExpStep.lean}{116}{Integral direction filter},
\src{src/Development/AutomorphicArithmetic/AutomorphicVarieties/ModularityArgument/ClassCounting/FProd.lean}{101}{cyclic residue cocycle},
\src{src/Development/AutomorphicArithmetic/AutomorphicVarieties/ModularityArgument/ClassCounting/STrans.lean}{411}{transfer over the projective lines}, and
\src{src/Development/AutomorphicArithmetic/AutomorphicVarieties/ModularityArgument/ClassCounting/LatticeBlockClasses/CJoin.lean}{5}{arithmetic cycle on the common lattice}.
